%%%%%%%%%%%%%%%%%%%%%%%%%%%%%%%%%%%%%%%%%%%%%%%%%%%%%%%%%%%%%%%%%%%%%%%%%%%%%%%%
%% 
\cleardoubleoddpage%  Make sure to start each chapter on a new odd page
\chapter{Conclusion and Outlook}
\label{chapter:conclusion}
In summary, this thesis introduces the fundamental notions underlying the scheduling problem addressed by the proposed \ac{IR} in \Cref{chapter:introduction}. Next, in \Cref{chapter:preliminaries-related-work} we present the language design, including justified choices made during its development. Furthermore, we introduce basic concepts of constrained optimization and representation of time in general. \Cref{chapter:methodology-language-design} demonstrates the design of our language and the reasoning behind certain decisions, highlighting the descriptive power of our language. After defining the language, we formalize the constraints mathematically in \Cref{chapter:formal-framework}, giving a clear definition of what a solution is and how to compare different solutions. In the process of defining the formal framework, we identify required properties, which serve as inputs for the solver. Subsequently, we establish the formal definition of a valid solution. Afterward, we define how to compare different solutions using cost functions, which can be used to find optimal solutions.

The main contribution of this thesis is the design and formalization of a highly expressive constraint language for personal scheduling. We have defined a set of constraint primitives that are significantly more expressive than they initially appear. As demonstrated in \Cref{chapter:methodology-language-design}, these primitives can be combined to express tasks with variable durations, conditional scheduling, and preference-based time windows. Crucially, we have provided rigorous formal semantics for all constraints in \Cref{chapter:formal-framework}, establishing a clear mathematical foundation for what constitutes a valid solution and how to compare different solutions through well-defined cost functions.

The defined \ac{IR} successfully addresses the goals outlined in \Cref{chapter:introduction} by enabling flexible scheduling across broad time spans while respecting complex temporal constraints. This language is designed to serve as an intermediate translation step between any high-level language and low-level solver-ready constraints. Both the high- and low-level languages are replaceable and not bound to any specific implementation, allowing for a wide range of applications as long as the high-level language can be compiled into the \ac{IR} and the low-level solver can handle the defined constraints.

As the proposed language only poses as an intermediate step, and it does neither handle solving the schedule itself nor provide a user interface to convert natural language input into \Ac{IR}. Hence, future work consists of first implementing a solver that can handle the defined constraints is of utmost importance. The defined formal framework serves as a basis for such an implementation. Such an implementation --- or rather its output --- can be checked with an existing implementation of a solution verifier\footnote{Repository with the implementation of a syntax compiler for the \Ac{IR} and a solution verifier: \url{https://github.com/h0p3zZ/BScThesis/tree/dev/_prog}}. Second, benchmarking different types of solvers with our defined constraints represents an interesting avenue for future research. Third, implementing a compiler from a high-level language into our \ac{IR} is beneficial to showcase the usability of our language. Finally, an intuitive way of defining tasks in a high-level language must also be provided. Such an input is most likely expressed in natural language, which can be parsed using \Ac{NLP} and/or \ac{LLM} techniques.

%% 
%%%%%%%%%%%%%%%%%%%%%%%%%%%%%%%%%%%%%%%%%%%%%%%%%%%%%%%%%%%%%%%%%%%%%%%%%%%%%%%%